\documentclass[10pt,a4paper]{amsart}
\usepackage{amsmath,amssymb,amsthm,graphicx}
\usepackage{verbatim}
\newenvironment{reference}{\comment}{\endcomment}
%\newenvironment{reference}{}{}
\newenvironment{slogan}{\comment}{\endcomment}
\newenvironment{history}{\comment}{\endcomment}

% For commutative diagrams we use Xy-pic
\usepackage[all]{xy}

% We use 2cell for 2-commutative diagrams.
\xyoption{2cell}
% Theorem environments.
%
\theoremstyle{plain}
\newtheorem{theorem}[subsection]{Theorem}
\newtheorem{proposition}[subsection]{Proposition}
\newtheorem{lemma}[subsection]{Lemma}

\theoremstyle{definition}
\newtheorem{definition}[subsection]{Definition}
\newtheorem{example}[subsection]{Example}
\newtheorem{exercise}[subsection]{Exercise}
\newtheorem{situation}[subsection]{Situation}

\theoremstyle{remark}
\newtheorem{remark}[subsection]{Remark}
\newtheorem{remarks}[subsection]{Remarks}

\numberwithin{equation}{subsection}

% Macros
%
\def\lim{\mathop{\mathrm{lim}}\nolimits}
\def\colim{\mathop{\mathrm{colim}}\nolimits}
\def\Spec{\mathop{\mathrm{Spec}}}
\def\Hom{\mathop{\mathrm{Hom}}\nolimits}
\def\Ext{\mathop{\mathrm{Ext}}\nolimits}
\def\SheafHom{\mathop{\mathcal{H}\!\mathit{om}}\nolimits}
\def\SheafExt{\mathop{\mathcal{E}\!\mathit{xt}}\nolimits}
\def\Sch{\mathit{Sch}}
\def\Mor{\mathop{\mathrm{Mor}}\nolimits}
\def\Ob{\mathop{\mathrm{Ob}}\nolimits}
\def\Sh{\mathop{\mathit{Sh}}\nolimits}
\def\NL{\mathop{N\!L}\nolimits}
\def\CH{\mathop{\mathrm{CH}}\nolimits}
\def\proetale{{pro\text{-}\acute{e}tale}}
\def\etale{{\acute{e}tale}}
\def\QCoh{\mathit{QCoh}}
\def\Ker{\mathop{\mathrm{Ker}}}
\def\Im{\mathop{\mathrm{Im}}}
\def\Coker{\mathop{\mathrm{Coker}}}
\def\Coim{\mathop{\mathrm{Coim}}}

% Boxtimes
%
\DeclareMathSymbol{\boxtimes}{\mathbin}{AMSa}{"02}

%
% Macros for moduli stacks/spaces
%
\def\QCohstack{\mathcal{QC}\!\mathit{oh}}
\def\Cohstack{\mathcal{C}\!\mathit{oh}}
\def\Spacesstack{\mathcal{S}\!\mathit{paces}}
\def\Quotfunctor{\mathrm{Quot}}
\def\Hilbfunctor{\mathrm{Hilb}}
\def\Curvesstack{\mathcal{C}\!\mathit{urves}}
\def\Polarizedstack{\mathcal{P}\!\mathit{olarized}}
\def\Complexesstack{\mathcal{C}\!\mathit{omplexes}}
% \Pic is the operator that assigns to X its picard group, usage \Pic(X)
% \Picardstack_{X/B} denotes the Picard stack of X over B
% \Picardfunctor_{X/B} denotes the Picard functor of X over B
\def\Pic{\mathop{\mathrm{Pic}}\nolimits}
\def\Picardstack{\mathcal{P}\!\mathit{ic}}
\def\Picardfunctor{\mathrm{Pic}}
\def\Deformationcategory{\mathcal{D}\!\mathit{ef}}
\usepackage{cleveref}
\crefname{originallemma}{original lemma}{original lemmas}
\Crefname{originallemma}{Original lemma}{Original lemmas}
\makeatletter
\newlabel{algebra-lemma-strictly-henselian-functorial-improve}{{08HV}{source}{}{sourcepointer-1}{}}
\newlabel{algebra-lemma-strictly-henselian-functorial-improve@cref}{{[originallemma][0][]08HV}{[1][0][]source}{}{}{}}
\makeatother
\crefname{originalreference}{original reference}{original references}
\Crefname{originalreference}{Original reference}{Original references}
\makeatletter
\newlabel{definition-tor-independent}{{0660}{source}{}{sourcepointer-2}{}}
\newlabel{definition-tor-independent@cref}{{[originalreference][0][]0660}{[1][0][]source}{}{}{}}
\makeatother
\crefname{originallemma}{original lemma}{original lemmas}
\Crefname{originallemma}{Original lemma}{Original lemmas}
\makeatletter
\newlabel{more-algebra-lemma-dumb-properties-henselization}{{07QM}{source}{}{sourcepointer-3}{}}
\newlabel{more-algebra-lemma-dumb-properties-henselization@cref}{{[originallemma][0][]07QM}{[1][0][]source}{}{}{}}
\makeatother
\crefname{originallemma}{original lemma}{original lemmas}
\Crefname{originallemma}{Original lemma}{Original lemmas}
\makeatletter
\newlabel{more-algebra-lemma-tor-independent}{{08HX}{source}{}{sourcepointer-4}{}}
\newlabel{more-algebra-lemma-tor-independent@cref}{{[originallemma][0][]08HX}{[1][0][]source}{}{}{}}
\makeatother
\crefname{originallemma}{original lemma}{original lemmas}
\Crefname{originallemma}{Original lemma}{Original lemmas}
\makeatletter
\newlabel{more-algebra-lemma-tor-independent-flat}{{08HW}{source}{}{sourcepointer-5}{}}
\newlabel{more-algebra-lemma-tor-independent-flat@cref}{{[originallemma][0][]08HW}{[1][0][]source}{}{}{}}
\makeatother
\crefname{originaldefinition}{original definition}{original definitions}
\Crefname{originaldefinition}{Original definition}{Original definitions}
\makeatletter
\newlabel{perfect-definition-tor-independent}{{08IA}{source}{}{sourcepointer-6}{}}
\newlabel{perfect-definition-tor-independent@cref}{{[originaldefinition][0][]08IA}{[1][0][]source}{}{}{}}
\makeatother
\crefname{originallemma}{original lemma}{original lemmas}
\Crefname{originallemma}{Original lemma}{Original lemmas}
\makeatletter
\newlabel{spaces-properties-lemma-describe-etale-local-ring}{{04KF}{source}{}{sourcepointer-7}{}}
\newlabel{spaces-properties-lemma-describe-etale-local-ring@cref}{{[originallemma][0][]04KF}{[1][0][]source}{}{}{}}
\makeatother
\crefname{originallemma}{original lemma}{original lemmas}
\Crefname{originallemma}{Original lemma}{Original lemmas}
\makeatletter
\newlabel{spaces-properties-lemma-geometric-lift-to-usual}{{05VN}{source}{}{sourcepointer-8}{}}
\newlabel{spaces-properties-lemma-geometric-lift-to-usual@cref}{{[originallemma][0][]05VN}{[1][0][]source}{}{}{}}
\makeatother
\title{Tor independence of algebraic spaces is etale local and agrees with the scheme notion}
\author{The Stacks Project Authors}
\date{stacks-project:spaces-perfect.tex:local-snapshot-20261012; GNU FDL 1.2 or later}
\begin{document}
\maketitle
\noindent\textit{Editorial note.} All delivered mathematical text is the author's own, copied from the pinned Stacks Project chapter file spaces-perfect.tex; no statement or proof step was rewritten, reordered or removed.\par\medskip
\noindent\textit{Editorial note.} The counted proof is the author's complete argument: etale-local characterizations via geometric point lifting, a strict henselization comparison with an explicit Tor base-change identity, and an affine converse using flatness and localized Tor vanishing, all carried out inside the delivered span\par\medskip
\noindent\textit{Editorial note.} The label lemma-tor-independent names the counted statement, and the proof environment immediately following it in the source is delivered verbatim as its proof.\par\medskip
\noindent\textit{Editorial note.} The counted statement is carried by the author's own lemma environment in the chapter source.\par\medskip
\noindent\textit{Editorial note.} The paragraph handling item (4) opens with the author's sentence 'it is clear that (4) comes down to the statement that Definition \textbackslash{}ref\{definition-tor-independent\} agrees with Derived Categories of Schemes, Definition \textbackslash{}ref\{perfect-definition-tor-independent\} when X, Y, and B are schemes' and then performs the strict-henselization comparison concretely; the load-bearing steps are inside the delivered span.\par\medskip
\noindent\textit{Allowed standard prerequisites (editorial).}\begin{itemize}
\item Tor independence of schemes as defined in Derived Categories of Schemes, and its extension to algebraic spaces via the strict-henselian local ring definition used in the source
\item Etale morphisms induce isomorphisms on strict local rings, and the functorial strict henselization description of etale local rings
\item Flat base change of Tor modules under henselization from More on Algebra
\end{itemize}\medskip
\ifdefined\hypertarget\hypertarget{sourcepointer-1}{}\fi
\noindent\textit{Source reference.} Original lemma 08HV: sources/stacks-chow/tags.txt, source lines 10793--10793. Prerequisite pointer invoked at proof line(s) 5025; it names a standard fact of a sibling chapter that the counted argument uses, and every step of the counted proof built on it is retained verbatim in the delivered body.\par\medskip
\ifdefined\hypertarget\hypertarget{sourcepointer-2}{}\fi
\noindent\textit{Source reference.} Original reference 0660: sources/stacks-spaces-perfect/pinned-original-source.tex, source lines 4942--4942. Prerequisite pointer invoked at proof line(s) 4996; it names a standard fact of the same chapter that the counted argument uses, and every step of the counted proof built on it is retained verbatim in the delivered body.\par\medskip
\ifdefined\hypertarget\hypertarget{sourcepointer-3}{}\fi
\noindent\textit{Source reference.} Original lemma 07QM: sources/stacks-chow/tags.txt, source lines 9837--9837. Prerequisite pointer invoked at proof line(s) 5035; it names a standard fact of a sibling chapter that the counted argument uses, and every step of the counted proof built on it is retained verbatim in the delivered body.\par\medskip
\ifdefined\hypertarget\hypertarget{sourcepointer-4}{}\fi
\noindent\textit{Source reference.} Original lemma 08HX: sources/stacks-chow/tags.txt, source lines 10795--10795. Prerequisite pointer invoked at proof line(s) 5071; it names a standard fact of a sibling chapter that the counted argument uses, and every step of the counted proof built on it is retained verbatim in the delivered body.\par\medskip
\ifdefined\hypertarget\hypertarget{sourcepointer-5}{}\fi
\noindent\textit{Source reference.} Original lemma 08HW: sources/stacks-chow/tags.txt, source lines 10794--10794. Prerequisite pointer invoked at proof line(s) 5036; it names a standard fact of a sibling chapter that the counted argument uses, and every step of the counted proof built on it is retained verbatim in the delivered body.\par\medskip
\ifdefined\hypertarget\hypertarget{sourcepointer-6}{}\fi
\noindent\textit{Source reference.} Original definition 08IA: sources/stacks-chow/tags.txt, source lines 10808--10808. Prerequisite pointer invoked at proof line(s) 4985; it names a standard fact of a sibling chapter that the counted argument uses, and every step of the counted proof built on it is retained verbatim in the delivered body.\par\medskip
\ifdefined\hypertarget\hypertarget{sourcepointer-7}{}\fi
\noindent\textit{Source reference.} Original lemma 04KF: sources/stacks-chow/tags.txt, source lines 5971--5971. Prerequisite pointer invoked at proof line(s) 5023; it names a standard fact of a sibling chapter that the counted argument uses, and every step of the counted proof built on it is retained verbatim in the delivered body.\par\medskip
\ifdefined\hypertarget\hypertarget{sourcepointer-8}{}\fi
\noindent\textit{Source reference.} Original lemma 05VN: sources/stacks-chow/tags.txt, source lines 7557--7557. Prerequisite pointer invoked at proof line(s) 5003; it names a standard fact of a sibling chapter that the counted argument uses, and every step of the counted proof built on it is retained verbatim in the delivered body.\par\medskip
\begin{lemma}
\label{lemma-tor-independent}
Let $S$ be a scheme. Let $B$ be an algebraic space over $S$.
Let $X$, $Y$ be algebraic spaces over $B$. The following are equivalent
\begin{enumerate}
\item $X$ and $Y$ are Tor independent over $B$,
\item for every commutative diagram
$$
\xymatrix{
U \ar[d] \ar[r] & W \ar[d] & V \ar[d] \ar[l] \\
X \ar[r] & B & Y \ar[l]
}
$$
with \'etale vertical arrows $U$ and $V$ are Tor independent over $W$,
\item for some commutative diagram as in (2) with (a) $W \to B$ \'etale
surjective, (b) $U \to X \times_B W$ \'etale surjective, (c)
$V \to Y \times_B W$ \'etale surjective, the spaces $U$ and $V$ are Tor
independent over $W$, and
\item for some commutative diagram as in (3) with $U$, $V$, $W$ schemes,
the schemes $U$ and $V$ are Tor independent over $W$ in the sense of
Derived Categories of Schemes, Definition
\ref{perfect-definition-tor-independent}.
\end{enumerate}
\end{lemma}

\begin{proof}
For an \'etale morphism $\varphi : U \to X$ of algebraic spaces
and geometric point $\overline{u}$ the map of local rings
$\mathcal{O}_{X, \varphi(\overline{u})} \to \mathcal{O}_{U, \overline{u}}$
is an isomorphism. Hence the equivalence of (1) and (2) follows.
So does the implication (1) $\Rightarrow$ (3). Assume (3) and
pick a diagram of geometric points as in
Definition \ref{definition-tor-independent}.
The assumptions imply that we can first lift $\overline{b}$ to a geometric
point $\overline{w}$ of $W$, then lift the geometric point
$(\overline{x}, \overline{b})$ to a geometric point $\overline{u}$ of $U$,
and finally lift the geometric point
$(\overline{y}, \overline{b})$ to a geometric point $\overline{v}$ of $V$.
Use Properties of Spaces, Lemma
\ref{spaces-properties-lemma-geometric-lift-to-usual}
to find the lifts.
Using the remark on local rings above we conclude that the condition
of the definition is satisfied for the given diagram.

\medskip\noindent
Having made these initial points, it is clear that (4) comes down to the
statement that
Definition \ref{definition-tor-independent}
agrees with
Derived Categories of Schemes, Definition
\ref{perfect-definition-tor-independent}
when $X$, $Y$, and $B$ are schemes.

\medskip\noindent
Let $\overline{x}, \overline{b}, \overline{y}$ be as in
Definition \ref{definition-tor-independent} lying over the points
$x, y, b$. Recall that
$\mathcal{O}_{X, \overline{x}} = \mathcal{O}_{X, x}^{sh}$
(Properties of Spaces, Lemma
\ref{spaces-properties-lemma-describe-etale-local-ring}) and similarly
for the other two. By Algebra, Lemma
\ref{algebra-lemma-strictly-henselian-functorial-improve} we see that
$\mathcal{O}_{X, \overline{x}}$ is a strict henselization of
$\mathcal{O}_{X, x} \otimes_{\mathcal{O}_{B, b}} \mathcal{O}_{B, \overline{b}}$.
In particular, the ring map
$$
\mathcal{O}_{X, x} \otimes_{\mathcal{O}_{B, b}} \mathcal{O}_{B, \overline{b}}
\longrightarrow
\mathcal{O}_{X, \overline{x}}
$$
is flat (More on Algebra, Lemma
\ref{more-algebra-lemma-dumb-properties-henselization}). By
More on Algebra, Lemma \ref{more-algebra-lemma-tor-independent-flat}
we see that
$$
\text{Tor}_i^{\mathcal{O}_{B, b}}(\mathcal{O}_{X, x}, \mathcal{O}_{Y, y})
\otimes_{\mathcal{O}_{X, x} \otimes_{\mathcal{O}_{B, b}} \mathcal{O}_{Y, y}}
(\mathcal{O}_{X, \overline{x}} \otimes_{\mathcal{O}_{B, \overline{b}}}
\mathcal{O}_{Y, \overline y})
=
\text{Tor}_i^{\mathcal{O}_{B, \overline{b}}}(
\mathcal{O}_{X, \overline{x}}, \mathcal{O}_{Y, \overline{y}})
$$
Hence it follows that if $X$ and $Y$ are Tor independent over $B$
as schemes, then $X$ and $Y$ are Tor independent as algebraic spaces over $B$.

\medskip\noindent
For the converse, we may assume $X$, $Y$, and $B$ are affine.
Observe that the ring map
$$
\mathcal{O}_{X, x} \otimes_{\mathcal{O}_{B, b}} \mathcal{O}_{Y, y}
\longrightarrow
\mathcal{O}_{X, \overline{x}} \otimes_{\mathcal{O}_{B, \overline{b}}}
\mathcal{O}_{Y, \overline y}
$$
is flat by the observations given above. Moreover, the image of the map
on spectra includes all primes
$\mathfrak s \subset
\mathcal{O}_{X, x} \otimes_{\mathcal{O}_{B, b}} \mathcal{O}_{Y, y}$
lying over $\mathfrak m_x$ and $\mathfrak m_y$.
Hence from this and the displayed formula of Tor's above we see that if
$X$ and $Y$ are Tor independent over $B$ as algebraic spaces, then
$$
\text{Tor}_i^{\mathcal{O}_{B, b}}
(\mathcal{O}_{X, x}, \mathcal{O}_{Y, y})_\mathfrak s = 0
$$
for all $i > 0$ and all $\mathfrak s$ as above. By
More on Algebra, Lemma \ref{more-algebra-lemma-tor-independent}
applied to the ring maps
$\Gamma(B, \mathcal{O}_B) \to \Gamma(X, \mathcal{O}_X)$
and
$\Gamma(B, \mathcal{O}_B) \to \Gamma(X, \mathcal{O}_X)$
this implies that $X$ and $Y$ are Tor independent over $B$.
\end{proof}

\end{document}
